% GENERATED FILE. Edit canonical lessons, then run npm run generate:books.
% canonical-source-hash: fnv1a64:4120d12f492c457a
% canonical-lessons: ZH-C181-beizi, ZH-C181-zhentou, ZH-C181-jian, ZH-W181-jian-room, ZH-C181-yushi, ZH-C181-chuanglian, ZH-R181-review-176-181

\chapter{Bed and Bath}
\label{ch:zh-bed-and-bath}
\hlchaptermodality{181}

\begin{chapteropening}
\textbf{By the end of this chapter:} \emph{I can name a quilt, a pillow, a bathroom and curtains, and count the rooms.}

Complete the last lesson of chapter 181: I can name a quilt, a pillow, a bathroom and curtains, and count the rooms.
\end{chapteropening}

\section[bèizi]{bèizi — quilt; duvet}

\label{lesson:ZH-C181-beizi}



\begin{quote}
Before the new one: say the Mandarin for to come over, then the Mandarin for to make a date with.
\end{quote}

\subsection*{You'll want to know: bèizi}
\emph{bèizi} — \textquotedblleft{}quilt; duvet\textquotedblright{}.

\begin{quote}
\emph{tiānqì lěng le wǒmen yào huàn bèizi} — It's turned cold: we need to change the quilt.
\end{quote}

Its characters come later; for now it is a word for the ear.

\begin{cousinweb}
\emph{bèizi} is built from \emph{bèi} 'bedclothes' and the noun ending \zh{子}, both inherited from Middle Chinese, \emph{bjeX} and \emph{tsiX} (Guangyun, in Baxter's transcription). In the word the second loses its tone.
\end{cousinweb}

\begin{grammarlens}[title={What you've built}]
One more word for everyday things.
\end{grammarlens}

\subsection*{Your turn}
\begin{itemize}
\raggedright
  \item \emph{Say it:} \emph{bèizi}
  \item \emph{Say it:} \emph{bèizi}, once more
  \item \emph{Say it:} \emph{tiānqì lěng le wǒmen yào huàn bèizi} — It's turned cold: we need to change the quilt.
  \item \emph{Recall:} say the Mandarin for to come over, then the Mandarin for to make a date with, then say \emph{bèizi} again
\end{itemize}

\begin{tcolorbox}[breakable,colback=teal!4,colframe=teal!35!black,title={Before you move on}]
What does bèizi mean? (\textquotedblleft{}Quilt; duvet\textquotedblright{}.) Say it once more.
\end{tcolorbox}
\section[zhěntou]{zhěntou — pillow}

\label{lesson:ZH-C181-zhentou}



\begin{quote}
Before the new one: say the Mandarin for to make a date with, then the Mandarin for a quilt.
\end{quote}

\subsection*{You'll want to know: zhěntou}
\emph{zhěntou} — \textquotedblleft{}pillow\textquotedblright{}.

\begin{quote}
\emph{zhège zhěntou tài gāo le} — This pillow is too high.
\end{quote}

Its characters come later; for now it is a word for the ear.

\begin{cousinweb}
\emph{zhěntou} is built from \emph{zhěn} 'pillow' and \emph{tóu} 'head', both inherited from Middle Chinese, \emph{tsyimX} and \emph{duw} (Guangyun, in Baxter's transcription); under its falling-tone reading the dictionary glosses the first with this very word. In the word the second loses its tone.
\end{cousinweb}

\begin{grammarlens}[title={What you've built}]
One more word for everyday things.
\end{grammarlens}

\subsection*{Your turn}
\begin{itemize}
\raggedright
  \item \emph{Say it:} \emph{zhěntou}
  \item \emph{Say it:} \emph{zhěntou}, once more
  \item \emph{Say it:} \emph{zhège zhěntou tài gāo le} — This pillow is too high.
  \item \emph{Recall:} say the Mandarin for to make a date with, then the Mandarin for a quilt, then say \emph{zhěntou} again
\end{itemize}

\begin{tcolorbox}[breakable,colback=teal!4,colframe=teal!35!black,title={Before you move on}]
What does zhěntou mean? (\textquotedblleft{}Pillow\textquotedblright{}.) Say it once more.
\end{tcolorbox}
\section[\texorpdfstring{\zh{间} (jiān)}{间 (jiān)}]{\zh{间} (jiān) — measure word for rooms}

\label{lesson:ZH-C181-jian}



\begin{quote}
Before the new one: say the Mandarin for a quilt, then the Mandarin for a pillow.
\end{quote}

\subsection*{You'll want to know: \zh{间}}
\textbf{\zh{间}} — \emph{jiān} — \textquotedblleft{}measure word for rooms\textquotedblright{}.

\begin{quote}
\textbf{\zh{后边有两间书房}} \emph{hòubian yǒu liǎng jiān shūfáng} — There are two studies at the back.
\end{quote}

The new character, in pieces: \zh{间} — seven strokes: the door piece \zh{门}, three strokes, and \zh{日} inside it, four strokes. The next lesson puts it in your hand.

\begin{cousinweb}
\textbf{\zh{间}} (\emph{jiān}) is inherited from Middle Chinese \emph{kean} (Guangyun, in Baxter's transcription), which writes it another way and glosses it 'a gap' and 'in between': the space between walls, and so a room.
\end{cousinweb}

\begin{grammarlens}[title={What you've built}]
One more word for everyday things.
\end{grammarlens}

\subsection*{Your turn}
\begin{itemize}
\raggedright
  \item \emph{Say it:} \emph{jiān}
  \item \emph{Say it:} \emph{jiān}, once more
  \item \emph{Say it:} \emph{hòubian yǒu liǎng jiān shūfáng} — There are two studies at the back.
  \item \emph{Recall:} say the Mandarin for a quilt, then the Mandarin for a pillow, then say \emph{jiān} again
\end{itemize}

\begin{tcolorbox}[breakable,colback=teal!4,colframe=teal!35!black,title={Before you move on}]
What does \zh{间} mean? (\textquotedblleft{}Measure word for rooms\textquotedblright{}.) Say it once more.
\end{tcolorbox}
\section[\texorpdfstring{\zh{间} (jiān)}{间 (jiān)}]{\zh{间} — \zh{门} with \zh{日} inside}

\label{lesson:ZH-W181-jian-room}



\begin{quote}
The word first, then the character inside it.

\begin{itemize}
\raggedright
  \item \emph{Recall:} say \emph{jiān} — \textbf{R1}, one lesson back, before it decays
\end{itemize}

\begin{quote}
\zh{门} \zh{时}
\end{quote}

\emph{Read it:} each one aloud, with nothing said first — \emph{mén}, \emph{shí}
\end{quote}

\subsection*{Script — \zh{门} with \zh{日} inside}
\hlblockfigure{\detokenize{figures/ZH-W181-jian-room-filmstrip.pdf}}{How \zh{间} is written, stroke by stroke}

\begin{quote}
\zh{间}
\end{quote}

\textbf{Two pieces, the frame first.} Count the pieces, not the strokes.

\textbf{\zh{门}, three strokes,} as you write it: a dot, a vertical down the left, and a horizontal that turns down the right side and hooks.

\textbf{Then \zh{日} inside it, four strokes,} as you write it, but smaller: a vertical, a horizontal that turns down, a bar across the middle, and the bottom bar.

3 + 4 = seven.

\subsection*{Your turn}
\begin{itemize}
\raggedright
  \item \emph{Write it:} \zh{间}, all seven, with the model in front of you
  \item \emph{Cover:} write it again from memory, and count the pieces, not the strokes
\end{itemize}

\begin{tcolorbox}[breakable,colback=teal!4,colframe=teal!35!black,title={Before you move on}]
How many strokes in \zh{间}? (\textbf{Seven}.) Which piece comes first? (\textbf{\zh{门}, the frame}.) What sits inside? (\textbf{\zh{日}}.)

Source: \href{https://raw.githubusercontent.com/chanind/hanzi-writer-data/68d10a4b21150cae5e1ebbd223eed289cf32d90c/data/\%E9\%97\%B4.json}{Hanzi Writer Data \zh{间}, pinned stroke paths and medians}.
\end{tcolorbox}
\section[yùshì]{yùshì — bathroom}

\label{lesson:ZH-C181-yushi}



\begin{quote}
Before the new one: say the Mandarin for a pillow, then the Mandarin for measure word for rooms.
\end{quote}

\subsection*{You'll want to know: yùshì}
\emph{yùshì} — \textquotedblleft{}bathroom\textquotedblright{}.

\begin{quote}
\emph{wǒmen jiā yǒu liǎng jiān yùshì} — We have two bathrooms.
\end{quote}

Its characters come later; for now it is a word for the ear.

\begin{cousinweb}
\emph{yùshì} is built from \emph{yù} 'to bathe' and \emph{shì} 'room', both inherited from Middle Chinese, \emph{yowk} and \emph{syit} (Guangyun, in Baxter's transcription).
\end{cousinweb}

\begin{grammarlens}[title={What you've built}]
One more word for everyday things.
\end{grammarlens}

\subsection*{Your turn}
\begin{itemize}
\raggedright
  \item \emph{Say it:} \emph{yùshì}
  \item \emph{Say it:} \emph{yùshì}, once more
  \item \emph{Say it:} \emph{wǒmen jiā yǒu liǎng jiān yùshì} — We have two bathrooms.
  \item \emph{Recall:} say the Mandarin for a pillow, then the Mandarin for measure word for rooms, then say \emph{yùshì} again
\end{itemize}

\begin{tcolorbox}[breakable,colback=teal!4,colframe=teal!35!black,title={Before you move on}]
What does yùshì mean? (\textquotedblleft{}Bathroom\textquotedblright{}.) Say it once more.
\end{tcolorbox}
\section[chuānglián]{chuānglián — curtain}

\label{lesson:ZH-C181-chuanglian}



\begin{quote}
Before the new one: say the Mandarin for measure word for rooms, then the Mandarin for a bathroom.
\end{quote}

\subsection*{You'll want to know: chuānglián}
\emph{chuānglián} — \textquotedblleft{}curtain\textquotedblright{}.

\begin{quote}
\emph{wòshì de chuānglián hěn hǎokàn} — The bedroom curtains are lovely.
\end{quote}

Its characters come later; for now it is a word for the ear.

\begin{cousinweb}
\emph{chuānglián} is built from \emph{chuāng} 'window' and \emph{lián} 'a hanging screen', both inherited from Middle Chinese, \emph{tsrhaewng} and \emph{ljem} (Guangyun, in Baxter's transcription), which glosses the second as a screen of thin strips.
\end{cousinweb}

\begin{grammarlens}[title={What you've built}]
One more word for everyday things.
\end{grammarlens}

\subsection*{Your turn}
\begin{itemize}
\raggedright
  \item \emph{Say it:} \emph{chuānglián}
  \item \emph{Say it:} \emph{chuānglián}, once more
  \item \emph{Say it:} \emph{wòshì de chuānglián hěn hǎokàn} — The bedroom curtains are lovely.
  \item \emph{Recall:} say the Mandarin for measure word for rooms, then the Mandarin for a bathroom, then say \emph{chuānglián} again
\end{itemize}

\begin{tcolorbox}[breakable,colback=teal!4,colframe=teal!35!black,title={Before you move on}]
What does chuānglián mean? (\textquotedblleft{}Curtain\textquotedblright{}.) Say it once more.
\end{tcolorbox}
\section[(dialogue)]{Looking back: from A Knock at the Door to Bed and Bath}

\label{lesson:ZH-R181-review-176-181}



\begin{quote}
Say the Mandarin for to knock at the door and a neighbour.
\end{quote}

\subsection*{Your turn}
\begin{itemize}
\raggedright
  \item \emph{Say it:} to knock at the door, a neighbour, to quarrel, to topple over, to catch fire
  \item \emph{Say it:} to recommend, tofu, full after eating, fresh, horrible to eat
  \item \emph{Say it:} normal, one's days, once upon a time, lately, in the small hours
  \item \emph{Say it:} a washing machine, washing powder, to hang out to dry, to iron, a wardrobe
  \item \emph{Say it:} to visit as a guest, the host, to call on someone, to come over, to make a date with
  \item \emph{Say it:} a quilt, a pillow, measure word for rooms, a bathroom, a curtain
  \item \emph{Read it:} \zh{常} \zh{间} — each one aloud, then a word that holds it
\end{itemize}

\begin{tcolorbox}[breakable,colback=teal!4,colframe=teal!35!black,title={Before you move on}]
To knock at the door? (\emph{qiāomén}.) A neighbour? (\emph{línjū}.) To quarrel? (\emph{chǎojià}.) To topple over? (\emph{dǎo}.) To catch fire? (\emph{zháohuǒ}.) To recommend? (\emph{tuījiàn}.) Tofu? (\emph{dòufu}.) Full after eating? (\emph{bǎo}.) Fresh? (\emph{xīnxiān}.) Horrible to eat? (\emph{nánchī}.) Normal? (\textbf{\zh{正常}}, \emph{zhèngcháng}.) One's days? (\textbf{\zh{日子}}, \emph{rìzi}.) Once upon a time? (\textbf{\zh{从前}}, \emph{cóngqián}.) Lately? (\emph{zuìjìn}.) In the small hours? (\emph{língchén}.) A washing machine? (\textbf{\zh{洗衣机}}, \emph{xǐyījī}.) Washing powder? (\emph{xǐyīfěn}.) To hang out to dry? (\emph{liàng}.) To iron? (\emph{yùn}.) A wardrobe? (\emph{yīguì}.) To visit as a guest? (\textbf{\zh{做客}}, \emph{zuòkè}.) The host? (\emph{zhǔrén}.) To call on someone? (\emph{kànwàng}.) To come over? (\textbf{\zh{过来}}, \emph{guòlai}.) To make a date with? (\emph{yuē}.) A quilt? (\emph{bèizi}.) A pillow? (\emph{zhěntou}.) Measure word for rooms? (\textbf{\zh{间}}, \emph{jiān}.) A bathroom? (\emph{yùshì}.) A curtain? (\emph{chuānglián}.)
\end{tcolorbox}
